\section{Discussion}\label{sec:discussion}
\paragraph{A grade with an explanation}
The practical contribution is a field-level assessment that retains the context in which information becomes sensitive. The marginal score identifies how much a field can add; its witness identifies the accompanying fields. This distinction explains why low single-field predictability need not imply low release risk. The same set-level evaluation supports a complementary output---minimal unsafe combinations---from which withholding strategies can be constructed. In the experiments, this combination-aware view substantially reduces the unsafe combinations left intact by single-field grading.

\paragraph{Efficient search over combinations}
The scan and the subset-specific readout address the main computational burden: evaluating many possible information sets. The ablations show that adapting the readout to each set is more consequential than increasing training effort per query. Random nonlinear features are already effective for average estimation error, while masked pre-training improves the selection of high-value backgrounds for certification. The all-column reconstruction experiment suggests a route to sharing this representation across changing target definitions. Inference graphs could likewise guide background selection by supplying physical relations~\cite{hu2026ignn}; the present evaluation focuses on exhaustive scans of small sets.

\paragraph{Scope of the assessment}
The reported risks depend on the attacker family, the available labelled history, and the sampled operating conditions. Random splits evaluate held-out observations from the same data collection; temporal changes may alter the measured values. The background budget also specifies the combinations covered by the audit. Although $K=2$ captures most measured marginal risk for nine targets, size-four combinations remain relevant to threshold-based withholding. Finally, field roles follow a common disclosure-based scenario rather than an assessment of each market's local rules, and the inference-graph baseline evaluates statistical propagation without expert labels or the attention training of~\cite{hu2026ignn}.

\section{Conclusion}\label{sec:conclusion}
This paper connects set-dependent inference leakage to field-level risk grading of power-system data. Budgeted marginal inference measures a field's largest predictive contribution over admissible backgrounds and identifies a witness combination. An amortized attacker makes the search practical through masked pre-training and a closed-form readout for each field set, while selected retrained backgrounds yield lower bounds on the marginal risk.

Across four data sets, single-field grading recovers only 13--17\% of critical fields. Scan--certify recovers up to 88\% with no false positives against the retrained reference; estimate-driven withholding leaves only 1--6\% of minimal unsafe combinations intact with near-minimal withholding. The results support assessing fields through their combinations, using subset-specific adaptation to make that assessment efficient, and selecting the background budget according to whether the audit prioritizes marginal-risk ranking or coverage of unsafe releases.

\section*{Data and code availability}
RTS-GMLC~\cite{barrows2020rts} and the AEMO NEMWEB archive~\cite{aemo2024nemweb} are public; the PJM and CAISO series were obtained from the operators' public data portals. Field-role tables, dispatch simulation, retraining, estimator, and analysis code will be made available upon publication.
